\chapter{Prior–Data Fitted Networks}


V této kapitole uvedeme samotnou architekturu Prior–Data Fitted Networks (dále jenom PFN).
Začneme krátkou motivací, kde popíšeme hlavní ideu této Depp Learning architektury.
Následně ukážeme mechanismy, které se schovávají uvnitř PFN, jinými slovy, jak PFN postupně zpracovává vstupní hodnoty, a to vše na příkladě Gaussovských procesů.

\section{Úvod do Prior–Data Fitted Networks}

Prior–Data Fitted Networks je Deep Learning (dále jenom DL) architektura, kterou v roce $2022$ uvedli Samuel Muller a Noah Hollmann v článku (REFERENCE).
Hlavní motivací pro vznik tohoto modelu je snaha posunout techniky Bayesovské statistiky na další úroveň spolu se strmým rozvojem v oblasti strojového učení. 
PFN je založen na proslulé architektuře transformer, která je viníkem dnešního pokroku všemožných Deep Learning modelu, a o niž víme díky velkým jazykovým modelům, a není to náhoda, existuje totiž jisté spojení mezi těmito architekturami. 
Ačkoli samotnou architekturu PFN a její technické nunace budeme podrobně diskutovat až v následující sekci, zde jenom uvedeme její základní charakteristiky.
PFN je tzv. encoder–only architektura (tj. neobsahuje dekóder blok, ale pouze enkóder blok transformeru), na kterou je pak napojena softmax hlava, pro vyhodnocování výsledků ve tvaru pravděpodobností.
Na vstup PFN dostává sekvenci dat.
V případě velkých jazykových modelu by se jednalo o slovy nebo věty, umístěné v jistém pořadí.
Tady je ta myšlenka shodná, ale namísto slov dostane síť neuspořádanou sekvenci dat, které získáme z prior rozdělení.  
Právě toto je důležitý intuitivní pojetí, na kterém je založená celá ta architektura.
Stejně jako po jazykovém modelu, který dostal sekvenci slov, budeme chtít aby odhadl celkový kontext a jako výsledek např. i následující neznáme slovo, které by mohlo následovat , budeme i po PFN chtít, aby po předložení sekvenci dat, zkusil uchopit celkový trend (neboli kontext) a předpovědět, co by mohlo následovat dále.
Což je velmi relevantní, když zpracováváme časové řady a chceme řešit regresní problém, jak tomu je např. během práce s Gaussovskými procesy.
\par 

Celý princip fungování teto metodiky lze popsat následujícím způsobem.
Jak již jsme diskutovali dříve, účelem Bayesovského přístupu ve statistice je využití apriorní znalosti nebo informace o jistém jevu v přírodě. 
A tento princip se dobře prolíná s obecnou úlohou aproximace funkcí pomocí neuronových sítí, protože při jejich návrhu vždy musíme enkódovat aprirorní znalosti o této funkce.
To se dělá buď přímo skrz architekturu konkrétního modelu nebo za pomocí různých regularizací. 
Každopádně je to velmi netriviální úloha, kterou je třeba vždy řešit než začneme řešit kteroukoliv úlohu pomocí DL nástrojů. 
Jako jedno z možných řešení se nabízí právě tzv. Bayesovská inference, jejíž hlavním předpokladem je, že se jistý jev nebo proces vyskytuje v reálném světě a disponujeme informacemi o jeho chování.
Tento předpoklad nám umožňuje definovat následující postup.    
Uvažujme jev nebo celý proces (to jest prior, pro jednoduchost si pod tím můžeme představit obyčejnou množinu dat), ze kterého můžeme neomezeně vzorkovat simulované datové sady $\mathcal{D}$ reprezentující různé prediktivní úlohy (body). 
Bez újmy na obecnosti předpokládejme, že tyto datové sady obsahují vstupní příznaky $\mathbf{x} \in \mathbb{R}^d$ a k nim příslušné výstupní třídy $y \in \mathcal{Y}$.
Během trénování sítě PFN náhodně vygenerujeme dávku (batch) obsahující $K$ takových datových sad. Pro každou $i$-tou sadu v dávce ($i \in \{1, \dots, K\}$) provedeme následující kroky:
\begin{enumerate}
    \item Zvolíme náhodně velikost kontextu $N_i$.
    \item Vygenerovanou datovou sadu rozdělíme na kontextovou (trénovací) část $\mathcal{D}^{(i)}_{train} = (\mathbf{x}_j^{(i)}, y_j^{(i)})_{j=1}^{N_i}$ a testovací část obsahující (pro jednoduchost) jeden testovací vzorek $(\mathbf{x}^{(i)}_{test}, y^{(i)}_{test})$. 
    Tedy předložíme modelu data ve tvaru $D^{(i)} = D^{(i)}_{train} \cup \{\mathbf{x}^{(i)}_{test}, y^{(i)}_{test}\}, \forall i \in K$.
\end{enumerate}

Následně natrénujeme náš model $q_\theta$ tak, že mu na vstup předložíme kontextovou sadu společně s testovacím příznakem a učíme ho predikovat správný výstup. K tomu minimalizujeme ztrátovou funkci Negative Log-Likelihood (NLL) přes celou dávku $K$ vzorků ve tvaru:

$$L(\theta) = - \sum_{i=1}^{K} \log q_{\theta}(y^{(i)}_{test} | \mathbf{x}^{(i)}_{test}, \mathcal{D}^{(i)}_{train})$$

Optimalizací této funkce (např. pomocí gradientního sestupu) hledáme parametry $\hat{\theta} = \arg\min_{\theta} L(\theta)$, čímž se síť učí provádět inferenci in-context na základě poskytnutých dat.
\par   

Toto schéma nám umožňuje aproximovat aposteriorní prediktivní rozdělení (PPD) pro libovolný problém, pro který jsme schopni definovat apriorní rozdělení (prior) a vzorkovat z něj data.
Zde je namístě si ujasnit terminologii.
\textbf{Aposteriorní rozdělení (Posterior)} popisuje naši nejistotu ohledně samotného generativního modelu $p$ (nebo jeho parametrů) poté, co jsme pozorovali trénovací data $\mathcal{D}_n$. 
Matematicky se jedná o podmíněné rozdělení modelu vzhledem k těmto datům, které se značí jako $\Pi(p|\mathcal{D}_n)$.
\textbf{Aposteriorní prediktivní rozdělení (PPD - Posterior Predictive Distribution)} naproti tomu popisuje nejistotu ohledně nových dat, které model dosud neviděl, jako je například predikce testovacího štítku $y$ pro nový vstupní příznak $x$, na základě již pozorovaných dat $\mathcal{D}_n$. 
Vypočítá se jako průměr všech možných podmíněných rozdělení $p(y|x)$ vážených jejich aposteriorní pravděpodobností. 
Matematicky je definováno pomocí integrálu:
$$ \pi(y|\mathbf{x},\mathcal{D}_n) = \int p(y|\mathbf{x})d\Pi(p|\mathcal{D}_n). $$
Takže účelem PFN je aproximovat právě PPD.

V praxi to znamená, že model je trénován na široké rodině úloh vygenerovaných z tohoto prioru. 
Tím se PFN naučí implicitně provádět bayesovskou inferenci nad daným prostorem funkcí.

V případě Gaussovských procesů to vypadá následovně: Jelikož známe přesný tvar rozdělení (např. Gaussovský proces s RBF kernelem), dokážeme z něj nekonečněkrát vzorkovat nová data. 
Navíc můžeme hyperparametrům RBF kernelu přiřadit vlastní rozdělení (tzv. hierarchický prior), čímž úlohu pro PFN zesložitíme.
V této chvíli se model musí naučit generalizovat (tj. najít obecné řešení) přes celou rodinu funkcí pro všemožné hodnoty hyperparametrů. 
Potom stačí takto natrénovanému modelu předložit několik známých bodů (tzv. kontext) a on nám v jediném průchodu sítí vrátí rozdělení, které chceme predikovat a snaží se co nejvíce aproximovat původní funkci.
\par

\section{Praktická motivace k PFN}
Abychom navázali na předchozí kapitolu, kde jsme uvedli dnes již standardní metody pro aproximaci PPD, ukážeme motivační příklady s porovnáním PFN s MCMC metody a variační inferencí.
Muller ve své práci tuto problematiku uvedl a na ní představil světu PFN architekturu (REFERENCEmuller).
Hlavní výhodou PFN je dle něj to, že na rozdíl od variační inference nebo metod Monte Carlo, že se učí aproximovat PPD přímo z datových vzorků.
Jediným a velmi důležitým předpokladem ale je, abychom mohli rychle a výpočetně levně vzorkovat z apriorního rozdělení. 
Mezitím MCMC požaduje nenormované posteriorní rozdělení.
Musíme umět jeho spočítat likelihood a také mít prior jako hustotu, ne jenom možnost vzorkovat. 
Pro variační inferenci potřebujeme znát sdruženou distribuci parametru a dat, často tyto distribuce jsou těžce spočitatelné.
Navíc často je optimalizace takových řešení výpočetně náročná (výpočet gradientů, ELBO optimalizace) a kvalita závisí na volbě variační rodiny rozdělení.
Dalším úskalím je, že pro každý nový dataset je nutné spustit celý proces znovu, což jenom zhoršuje časovou a výpočetní náročnost při řešení jistého problému.
Výhodou těchto klasických řešení je, že např. MCMC konverguje k pravému posterioru v limitě nekonečně mnoha vzorků.
Type II – MLE (speciální případ variační inference) konverguje k pravým hyperparametrům v limitě nekonečně dat za předpokladu korektně specifikovaného modelu. 
Navíc pro všechny modely platí, že systematická chyba takových aproximací, tzv. bias, jde k nule.
\par   

PFN obrací situaci.
Model se předtrénován offline na obrovské mase syntetických datasetů z prioru a inference se pak provádí jediným dopředním průchodem (forward passem) sítí. 
Strukturální požadavek je výrazně slabší – stačí umět vzorkovat z apriorní distribuce, není potřeba hustota ani její gradient. 
To otevírá třídu priorů, kterou klasické metody pojmout neumí, např. BNN s libovolnou architekturou nebo ručně psané generativní procedury. 
Inference je řádově rychlejší, protože model je předtrénovaný.
Müller ukazuje cca $200\times$ zrychlení proti type II – MLE a přibližně $1 000\times – 8 000\times$ proti NUTS, navíc na rozdíl od NUTS, PFN neprovádí žádné optimalizace před inferencí (tj. žádný warmup, žádná adaptace, žádné chain diagnostics). 
Konečně, PFN není trénován na rekonstrukci posterioru bod po bodu, ale přímo na predikční kvalitě skrz cross-entropy s PPD. 
Nevýhody jsou strukturální. 
Jak se ukáže dále, nemáme vždy zaručeno to, že bias bude klesat k nule. 
Naopak např. u architektury transformer bias nikdy nezmizí.
To je kvalitativně jiná situace než u NUTS nebo type II – MLE, které teoreticky konvergují přesně k prioru a jejich bias postupně mizí. 
Další nevýhodou je např. to, že pokus testovací dataset pochází z jiné distribuce než trénovací prior, PFN může selhat nebo minimálně výsledky nebudou tak přesné. 
NUTS je oproti tomu robustní v tom smyslu, že jeho výstupem je vždy korektní posterior pro právě ten model, který byl explicitně specifikován, bez ohledu na původ dat. 
Nakonec je třeba mít na paměti, že PFN musí být předem natrénovaný na konkretním problému a s obrovským množstvím dat.
Např. pokud bychom chtěli aproximovat Gaussovský proces, je třeba to PFN předem natrénovat na dostatečně velké rodině funkcí a pestrou distribucí parametrů/hyperparametrů.
Pokud by samozřejmě někdo to udělal za nás, a používali bychom předtrénovaný PFN jako hotový produkt (via dnešní LLM), pak bychom pro libovolný problém dostávali okamžitou aproximaci.

\section{Vnitřní architektura PFN}

V této sekci se pokusíme popsat procesy, které se dějí uvnitř PFN transformeru.
Naším účelem není poskytnout kompletní popis architektury transformer obecně, za co odpovídají jednotlivé komponenty nebo parametry.
Předpokládáme, že čtenář je již seznámen se základy architektury transformer.
Pokud tak tomu není, doporučujeme nastudovat alespoň základní pojmy a principy této architekury. 
Potřebné zájemce jenom odkážeme na příslušnou kapitolu v knize (REFERENCE UDNERSTANDING DEEP LEARNING).
Přesto pevně věříme, že našemu výkladu budou rozumět i laici v oblasti strojového učení, neboť zde nepotřebujeme pokročilé znalosti nebo složité pojmy a zkusíme vše popsat co nejjednodušeji.
Pokud se takové ve výkladu objeví, slouží jenom pro odborníci v této oblasti a na porozumění hlavních principu nebudou mít velký vliv. 
Dále opět podotýkáme, že učelem práce není rozebírat konkretní implementaci PFN a související se s tím nuance.
Navíc PFN je relativně mladá architektura, které se neustále vylepšuje a aktualizuje, proto pro nás nemá smysl jít do takových detailů.
Aktuální verzi zdrojového kód PFN naleznete zde (REFERENCEpfngit).
\par 


V předchozí sekci jsem již zmínili, že PFN je transformer, který se trénuje na obrovském množství synetických datasetů, které vzorkujeme z prioru.
Tímto způsobem se PFN naučí provádět Bayesovskou inferenci a po natrénování stačí jeden forward pass sítí a model vrátí kompletní PPD pro nové testovací body na vstupu. 
V této sekci budeme uvažovat, že naším priorem je Gaussovsky proces s RBF kenrlem.
Dále chceme upozornit, že celý proces, který se odehrává uvnitř PFN od vstupu po výstup, popíšeme s použitím konkrétní konfigurace, kterou jsme běžně používali během zkoumání PFN:
\begin{enumerate}
    \item Dimenze embidingu \textit{EMSIZE = 512}.
    \item Počet attention hlav v každé vrstvě \textit{NHEAD = 8}. 
    \item Dimenze vnitřních skrytých vrstev dopředných neuronových sítí \textit{NHID = 1024}.
    \item Počet vrstev \textit{NLAYERS = 6}.
\end{enumerate}
Jelikož se vše chystáme navíc ukázat pro případ $1D$ GP regrese, použijeme úplně klasické nastavení PFN, které uvádějí Muller a Hollmann v (REFERENCEmuller), konkrétně s těmito parametry:
\begin{itemize}
    \item \lstinline|features_per_group = 1| — nastaven na jedničku, protože řešíme $1D$ regresi.
    \item \lstinline|attention_between_features = False| — slouží pro zpracování tabulkových dat; pokud je zapnutý, jedná se o architekturu \textbf{TabPFN}, kterou v této práci neuvažujeme (viz REFERENCEtabpfn).
\end{itemize}
Nakonec uvažujme, že provádíme inferenci pro GP regresi s $20$ trénovacími a $50$ testovacími body. 
Teď už máme nachystané formální parametry a nastavení, což nám umožňuje se pustit na samotný popis.
Celý proces si rozdělíme na několik fází.

\subsection{Fáze 1 - Meta–Learning}
Pro obecnost ukážeme, jak probíhá učení pro případ, kdy hyperparametry pro RBF kernel nejsou známe a pochází z nějakého rozdělení. 
Případ fixních hyperparametrů je pouze speciálním případem tohoto problému a všechno, co napíšeme dále, platí i pro tento speciální případ. 
Meta–Learning je paradigma strojového učení, které spočívá v tom, že namísto toho, abychom model naučili řešit jeden konkrétní fixní problém, naučíme ho řešit celou třídu podobných úkolů.
Konkrétně v našem případě by to znamenalo, že model naučíme na základě dat z GP s RBF kernelem, kde mu ukážeme celé spektrum všemožných hyperparametrů, které se v RBF kernelu vyskytují. 
Potom stačí jenom předložit modelu těch $20$ kontext bodů a hned bude schopen aproximovat GP, ze kterého tyto body pocházejí. 
V předchozí sekci jsem již drobně popsali, jak funguje trénink PFN, formálněji a mnohem podrobněji lze tento proces zapsat takto:

\clearpage
\begin{algorithm}[!t]
\caption{Meta-trénink sítě PFN na syntetických GP datech}
\label{alg:pfn_training}
\begin{algorithmic}[1]
\Require Model $q_\theta$, apriorní rozdělení hyperparametrů $\Pi$, velikost kroku $\eta$, velikost dávky $m$
\While{není dosaženo konvergence}
    \For{$j = 1, \ldots, m$}
        \State Vzorkuj hyperparametry GP: $(\ell_j, \sigma_j, \sigma_{n,j}) \sim \Pi$
        \State Vzorkuj vstupní body $x_1, \ldots, x_N$ rovnoměrně z $[0, 1]$
        \State Vzorkuj výstupy z GP s daným kernelem: $(y_1, \ldots, y_N) \sim \mathcal{GP}(0,\, k_{\ell_j, \sigma_j}) + \mathcal{N}(0, \sigma_{n,j}^2 I)$
        \State Rozděl na kontextovou sadu $\mathcal{D}_j = \{(x_i, y_i)\}_{i=1}^{n}$ a testovací body $\{(x_i^*, y_i^*)\}_{i=1}^{t}$
    \EndFor
    \State Předej $q_\theta$ kontexty $\mathcal{D}_j$ a testovací vstupy $x_i^*$; model vrátí prediktivní distribuce přes biny
    \State Spočítej ztrátu (NLL):
    \[
        \mathcal{L}(\theta) = -\frac{1}{m}\sum_{j=1}^{m}\sum_{i=1}^{t} \log q_{\theta}(y_i^* \mid x_i^*, \mathcal{D}_j)
    \]
    \State Aktualizuj: $\theta \leftarrow \theta - \eta\,\nabla_\theta \mathcal{L}(\theta)$
\EndWhile
\State \Return $\hat{\theta}$
\end{algorithmic}
\end{algorithm}

Tímto způsobem se dostáváme do situace, kdy model vidí miliony různých datasetů generovných z GP prioru a učí se pravidla, dle kterých pak musí uhádnout, jak vypadá posteriorní predikce na základě předložených dat.
Tohle je velmi důležitá myšlenka, PFN se neučí aproximovat konkrétní funkce, učí se inferenční pravidlo.

\subsection{Fáze 2 - Vstup do modelu}
Během inference model dostane tři tenzory: 
\begin{enumerate}
    \item Souřadnice trénovacích bodů \texttt{x\_train}.
    \item Naměřené hodnoty na trénovacích bodech \texttt{y\_train}. 
    \item Souřadnice testovacích bodů \texttt{x\_test}.
\end{enumerate}

Jak už plyne z logiky věci, \texttt{x\_train} a \texttt{y\_train} jsou dva různé objekty, první říká, kde jsme měřili (souřadnice na ose x), a druhý říká, co jsme v těchto místech naměřili.
Tak dohromady tvoří datové body $(x_i, y_i)$.
Souřadnice \texttt{x\_test} nám říká, kde chceme ty predikce, přitom model by měl vrátit pravděpodobnostní rozdělení pro příslušné y–hodnoty, ale k tomu se dostaneme později v naším výkladu.
Následně se \texttt{x\_train} a \texttt{x\_test} spojí do jedné sekvenci a model s nimi pracuje jako s jedním objektem.
Přičemž modelu současně explicitně řekneme, že prvním $20$ souřadnic jsou trénovací hodnoty a zbytek jsou testovací ($50$ testovacích souřadnic), jak jsme si to definovali na začátku sekce.
Pro y–hodnoty je ta situace podobná, vytvoříme jednou sekvenci pro \texttt{y\_train} a \texttt{y\_test} hodnoty, až na to, že místo test hodnot vložíme \texttt{NaN} hodnoty, protože právě ty musí model odhadnout a nechceme je do modelu vkládat v žádné podobě.
Pro zájemce uvedeme, jak to pak vypadá v kódu.
V metodě \texttt{forward()} se \texttt{x\_train} a \texttt{x\_test} spojí pomocí funkce konkatenace:
\begin{lstlisting}[language=Python]
x = torch.cat((x_train, x_test), dim=1)  # (1, 70, 1)
\end{lstlisting}
Jedná se o tenzor, kde první dimenze je počet datasetů, které zpracováváme naraz, druhá dimenze je počet bodů v sekvenci a třetí dimenze je počet příznaků, který jsme na začátku této sekce nastavili na $1$.
Pro y–hodnoty analogicky:
\begin{lstlisting}[language=Python]
y = torch.cat((y_train, NaN * torch.zeros(1, 50, 1)), dim=1)  # (1, 70, 1)
\end{lstlisting}
Zde je explicitně vidět, že test hodnoty se nahradí za \lstinline|NaN|.
Dále si myslíme, že je namístě říct jedno upřesnění, aby došlo k nedorozumění. 
To, že modelu v y–hodnotách předkládáme \texttt{NaN} neznamená, že pak budeme chtít tyto hodnoty nahradit za odhady někde v průběhu zpracovávání, je to čistě technická záležitost. 
Zůstanou \texttt{NaN} po celou dobu.
Výsledky obdržíme v jiné podobě, opět, budeme o tom diskutovat dále.
\par   

Dále následuje tzv. \textit{Feature grouping}, o parametru, který určuje počet příznaku jsme také zmiňovali na začátku sekce.
Připomínáme, že pro náš jednoduchý případ $1$D regrese je nataven na $1$.
Tento mechanismus slouží hlavně pro případ, kdy chceme zpracovávat tabulkové hodnoty, anebo pracovat s vícerozměrnými daty, ale to se týká hlavně verze modelu \textbf{TabPFN}.
V tabulkových datech s 12 sloupci (věk, plat, výška, \dots) model potřebuje tyto sloupce nějak zpracovat. 
Parametr \texttt{features\_per\_group} určuje, kolik sloupců se seskupí do jednoho tokenů.
Pro náš $1$D GP je grouping triviální: máme $1$ přiznak, tedy $1$ grupu. 
Tenzor $x$ se tak přeformátuje z $(1, 70, 1)$ na $(1, 70, 1, 1)$.
\par   

Dalším krokem je převod dat na vstupu na tzv. \textit{embeding}, tj. do řeči, které rozumí transformer. 
Ve skutečnosti se jedná o vytváření nového prostoru, ve kterém pak PFN bude operovat s daty, a kterému rozumí víc, než kdyby musel pracovat pouze se sekvenci dat.


\textbf{X–encoder}:

Zde to je jednoduchá lineární vrstva, která vezme jednu konkretní číselnou hodnotu $x$–souřadnice a vynásobí ji maticí vah $1\times512$ (parametr velikosti embedingu jsme nastavili nahoře).
Výsledkem je $512$–dimenzionální vektor (\textit{embeding}).
Tímto způsobem jsme zakódovali polohu bodu v tom novém prostoru, ze kterého model umí jednoduše číst a operovat v něm.
Výstupem je \texttt{embedded\_x} tvaru $(1, 70, 1, 512)$

\textbf{Y–encoder}:

Pro y–hodnoty je ta situace zajímavější, jelikož v předchozím kroku, jsme si řekli, že pro testovací hodnoty modelu předložíme \texttt{NaN} hodnoty. 
A teď otázkou je, jak tyto hodnoty zakódovat do \textit{embedingu}.
K tomu slouží metoda \lstinline|NanHandlingEncoderStep|, která nejdříve transformuje skalární y–hodnotu na vektor délky 2:

\begin{enumerate}
    \item Pro trénovací pozice ($y$ existuje): $[y, -1]$ — skutečná hodnota a indikátor "data jsou k dispozici".
    \item Pro testovací pozice ($y =$ \texttt{NaN}): $[0, +1]$ — nulová hodnota a indikátor "data chybí". 
\end{enumerate}

Pak lineární vrstva opět převede dvourozměrný vektor do $512$–dimenzionálního prostoru.
Výstupem je \texttt{embedded\_y} tvaru $(1, 70, 512)$, tady je o jeden parametr méně, protože příznaky mezi skupinami se nevztahují na $y$–hodnoty.
\par   

V této chvíli máme dva embedingy, cílem ale je, mít jeden hromadný embdeding, který v sobě nese informace o souřadnici $x$ a odpovídající ji hodnotě $y$.
Zde opět existují dva způsoby implementace v PFN: pro \lstinline|attention_between_features=False| a \lstinline|attention_between_features=True|, kde poslední verze je opět určena pro práci s vícedimenzionálními daty. 
Pro tu první verzi je postup velmi jednoduchý – \texttt{embedded\_x} a \texttt{embedded\_y} jednoduše sečteme. 
Výsledkem je jeden token pro odpovídající jednu konkretní pozici, tj, tento token kombinuje informaci o poloze $x$ a hodnotě $y$ v jednom vektoru.
Sčítání tady funguje, protože oba embidingy jsou již stejného tvaru a jsou ve stejném $512$–dimenzionálního prostoru a model bude umět z tohoto součtu dekódovat obě informační položky díky předchozímu postupu.
V kódu to pak formálně vypadá následovně:
\begin{lstlisting}[language=Python]
embedded_input = embedded_x + embedded_y.unsqueeze(2)  
# (1, 70, 1, 512) + (1, 70, 1, 512) -> (1, 70, 1, 512)
\end{lstlisting}
Druhý režim zde již popisovat nebudeme, zájemce odkážeme na zdrojové kódy na GitHub, které jsme uvedli na začátku kapitoly.

\subsection{Fáze 3 – Transformer vrstvy}
Tenzor \texttt{embedded\_input} pak musí projít $6$ transformer vrstvami (jejich počet jsme též definovali na začátku sekce).
Jakmile embedingy vstupují do transformer vrstev je zvykem jim říkat token. 
To znamená, že objektům \texttt{embedded\_x} a \texttt{embedded\_y} bychom říkali \texttt{x\_token} a \texttt{y\_token}.
Ačkoliv jsme tyto dva embedingy sloučili do jednoho, pořad musíme na ně nahlížet jako na samostatné objekty, protože naším cílem je nakonec získat \texttt{y\_token}, který v sobě bude obsahovat predikce pro testovací hodnoty.
Sloučení se tedy provádí pouze pro technické potřeby architektury.
Vraťme se ale zpátky k věci.
Každá vrstva v naší konfiguraci provádí dvě hlavní operace:
\begin{enumerate}
    \item Self–attention přes datové body (neplést s attention between features)
    \item Feedforward MLP
\end{enumerate}

\textbf{Self–attention}:

Tenzor se transponuje na \lstinline|(1, 1, 70, 512)|, aby sekvence $70$ bodů byla na dimenzi \lstinline|dim=2| (standardní konvence pro attention).
Multi-head attention s $8$ hlavami pak počítá:

Pro každou hlavu $h$ (z $8$):
$$Q^{(h)} = x \cdot W_Q^{(h)}, \quad K^{(h)} = x \cdot W_K^{(h)}, \quad V^{(h)} = x \cdot W_V^{(h)}$$
kde $W_Q, W_K \in \mathbb{R}^{512 \times 64}$ a $W_V \in \mathbb{R}^{512 \times 64}$ (protože $512 / 8 = 64$).

Klíčový detail: $K$ a $V$ se počítají pouze z trénovacích bodů. V kódu:
\begin{lstlisting}[language=Python]
attention_src_x = x[:, :single_eval_pos].transpose(1, 2)  
# first 20 positions
\end{lstlisting}
Attention matice je pak:
$$A = \text{softmax}\left(\frac{Q \cdot K^T}{\sqrt{64}}\right)$$

Pro trénovací body (pozice $0$–$19$) je $Q$ i $K$ ze stejných $20$ bodů, takže matice je $20 \times 20$ — trénovací body se dívají na sebe navzájem.

Pro testovací body (pozice $20$–$69$) $Q$ pochází z testovacích bodů, ale $K$ pochází z trénovacích bodů. Matice je $50 \times 20$ — testovací body se dívají na trénovací body.

Tohle je přímá analogie ke GP predikci. Ve formuli $\mu_* = k(x_*, X) K^{-1} y$ se testovací body taky \uv{dívají} na trénovací body přes kernel $k(x_*, X)$. Rozdíl je, že GP používá explicitní RBF kernel, zatímco PFN se naučil vlastní implicitní \uv{kernel} zakódovaný ve váhách $W_Q$ a $W_K$.

Výstup attention:
$$\text{Output}^{(h)} = A \cdot V^{(h)}$$

Všech $8$ výstupů (každý $\in \mathbb{R}^{64}$) se zřetězí a projde výstupní projekcí $W_O$:
$$\text{MultiHead} = \text{Concat}(\text{Output}^{(1)}, \dots, \text{Output}^{(8)}) \cdot W_O$$

Po attention následuje reziduální spojení (vstup $+$ výstup attention) a LayerNorm — normalizace, která stabilizuje trénink.
Reziduální spojení je technika, která přidává původní vstupní reprezentaci (třeba embedding datového bodu) k výstupu nelineární vrstvy (například MLP) těsně předtím, než tento výsledek pošle dál do sítě.
Reziduální spojení zajišťuje, že model nezačíná optimalizaci chaotickým hádáním, nýbrž bezpečným průtokem dat, ke kterému postupně přidává potřebnou komplexitu.
Co se tyče LayerNorm, ten normalizuje embedding na nulový průměr a jednotkový rozptyl přes posledních $d$ dimenzí (v našem případě $512$). 
Pro vektor $x$ délky $512$:
$$\text{LayerNorm}(x) = \frac{x - \mu}{\sigma + \epsilon},$$
kde $\mu$ a $\sigma$ jsou průměr a směrodatná odchylka přes $512$ komponent vektoru a $\epsilon = 10^{-5}$ je malá konstanta pro numerickou stabilitu.

\textbf{MLP – Multi-Layer Perceptron}:

MLP je dvouvrstvá dopředná síť:
$$\text{vstup}\;(512) \;\to\; \text{Linear}(512 \to 1024) \;\to\; \text{GELU} \;\to\; \text{Linear}(1024 \to 512) \;\to\; \text{výstup}\;(512)$$
Vezme $512$-dimenzionální vektor, rozšíří ho do $1024$ dimenzí přes lineární vrstvu, aplikuje nelineární aktivaci GELU (hladká varianta ReLU), a zase zúží zpět na $512$.

GELU (Gaussian Error Linear Unit) je funkce
$$\text{GELU}(x) = x \cdot \Phi(x),$$
kde $\Phi$ je distribuční funkce standardního normálního rozdělení.
Pro velká kladná $x$ se chová jako identita, pro velká záporná $x$ vrací přibližně nulu.

Proč nakonec musíme přidávat MLP je pro zkušeného ve strojovém učení čtenáře poněkud triviální otázka, vysvětlíme i pro ostatní čtenáře.
Attention je totiž lineární operace (ačkoliv se to nezdá), počítá pouze vážené průměry.
MLP do tohoto procesu přidává nelinearitu, bez něj by transformer mohl dělat jenom lineární transformace s embedingy.
Avšak ve strojovém učení je to nedostačující, chceme zachytit a složitější transformace a operace, čímž zvyšujeme generalizační schopnosti modelu.
A právě díky této vlastnosti může PFN aproximovat např. inverzi matice $\mathbf{K}^{-1}$ v GP.
Váhy druhé lineární vrstvy jsou inicializovány na nulu.
To znamená, že na začátku tréninku MLP vrací nulový vektor, a tak se tato vrstva na začátku chová jako identita.
To se však mění v průběhu tréninku.
Model začíná od jednoduchého stavu a postupně se učí komplexnější transformace.
Po MLP opět následuje reziduální spojení a normalizace.
\par 

Toto se opakuje tolik, kolik vrstev obsahuje model, v našem případě $6$ krát.
Po každé vrstvě se embedding každého bodu obohacuje o informaci z ostatních bodů (přes attention) a transformuje (přes MLP). 
Po $6$ vrstvách embedding testovacího bodu obsahuje dostatečnou informaci k tomu, aby dekóder z něj vyrobil prediktivní distribuci.
\par   

Po průchodu $6$ vrstvami dostáváme opět nový tenzor.
Ten obsahuje zpracované reprezentace jak trénovacích, tak i testovacích bodů.
Naším úkolem teď by bylo extrahovat testovací hodnoty $y$ (neboli $y$–tokeny, jak jsme diskutovali výše), respektive jejich predikce, které model postupně vyrobil.
Zbytek je pro účely inference nepodstatný.
Platí totiž, že celá architektura PFN je postavena tak, že se právě v $y$-tokenu postupně akumulují informace přes všechny vrstvy transformeru přes attention z $x$–tokenů a z kontextu trénovacích bodů.
Na konci celého procesu $y$–token pro testovací pozici obsahuje $512$–dimenzionální vektor, který implicitně kóduje celou posteriorní prediktivní distribuci (PPD) pro daný trénovací bod.
Uveďme opět pro zájemce, jak tento proces vypadá uvnitř, což by mělo pomocí k pochopení, co se v tomto kroku odehrálo:
\begin{lstlisting}[language=Python]
# Split into train and test part
test_encoder_out = encoder_out[:, current_context_len:, :]
# (1, 50, 1, 512) -- test positions (20-69)

train_encoder_out = encoder_out[:, :current_context_len, :]
# (1, 20, 1, 512) -- train positions (0-19)
\end{lstlisting}

Proměnná \lstinline|current_context_len| je rovna \lstinline|single_eval_pos = 20|.

V dalším kroku se extrahuje $y$-token — poslední token v dimenzi feature groups:
\begin{lstlisting}[language=Python]
test_encoder_out = test_encoder_out[:, :, -1, :]
# (1, 50, 512) -- only y-token for test positions
\end{lstlisting}

\subsection{Fáze 4 – Dekóder, převod embdedingu na logity.}
Jsme si již řekli, že PFN je encoder–only architektura, tj. decoder blok, který je součastí obecné architektury transformer, v něm chybí.
Decoder blok ale je napojen na hlavu, která v sobě obsahuje opět MLP vrstvu a softmax vrstvu.
Tyto součástky ma ovšem potřebovat budeme.
Na konci předchozí fáze jsme získali tenzor \texttt{test\_encoder\_out} o tvaru $(1, 50, 512)$, který v sobě obsahuje pro každý z $50$ testovacích bodů jeden $512$–dimenzionální embeding.
Teď potřebujeme z tohoto vektoru vytvořit prediktivní distribuci.
To se opět udělá pomocí MLP (jednoduchá dopředná síť), která je doplněná o jednu novinku.
Celkově schéma vypadá následovně:
$$\text{vstup}\;(512) \;\to\; \text{Linear}(512 \to 1024) \;\to\; \text{GELU} \;\to\; \text{Linear}(1024 \to \texttt{num\_bars}) \;\to\; \text{výstup}\;(\texttt{num\_bars})$$
Jednotlivé prvky v schématu, který jsme uvedli na začátku sekce, plní následující funkci:
\begin{enumerate}
    \item \lstinline|Linear(512 -> 1024)|: Vezme $512$-dimenzionální embedding a promítne ho do $1024$ dimenzí. To je lineární transformace — násobení maticí vah $W_1 \in \mathbb{R}^{512 \times 1024}$.
    \item \lstinline|GELU()|: Aplikuje nelineární aktivaci. Bez ní by dvě lineární vrstvy za sebou byly ekvivalentní jedné lineární vrstvě (součin matic je matice). GELU umožňuje decoderu zachytit nelineární vztahy mezi embeddingem a výstupní distribucí.
    \item \lstinline|Linear(1024 -> num_bars)|: Zúží z $1024$ dimenzí na \lstinline|num_bars| (řekněme $1000$). Výstupem je $1000$ čísel — \textbf{logity} — jedno pro každý bucket.
\end{enumerate}
Je vidět, že tady se objevuje nová proměnná \texttt{num\_bars}, jedná se totiž o analogii binu v histogramech, který se zavádí pro tzv. \textbf{BarDistribution}.

\textbf{BarDistribution}

BarDistribution je po–částech konstantní pravděpodobnostní rozdělení.
V článku (REFERENCEmuller) se mu říká \textit{Riemann distribution}, jedná se o diskretizaci spojité distribuce.
Vizuálně to vypadá takto:
Avšak je důležité odlišit tento koncept od klasického histogramu.
Histogram odhaduje hustotu z dat.
Hlavním principem je, že počítáme klik naměřených hodnot padlo do každého binu, a to určuje výšku jednotlivých sloupců.
BarDistribution je naopak výstup modelu, tj. my neděláme měření a nepočítáme, kolik hodnot padlo do toho či oného binu.
Právě samotný model určuje výšku sloupců dle nějakých svých pravidel.
Neuronová síť na výstupu vrátí $n$ čísel, každé číslo říká, jak vysoký má být jeden sloupec. 
To jest distribuce určuje síť, to umožňuje pak modelovat pestrou škálu různých rozdělení a nejsme omezení jenom normálním, exponenciálním rozdělením apod.
Binům se v BarDistribution říká \textit{bucket}.
V implementaci je BarDistribution definován pomocí hranic (parametr \texttt{borders}), tím definujeme interval, na kterém je tato funkce definována.
Může to vypadat například takto:\newline
\verb|borders = [-3.0, -2.994, -2.988, ..., 0.0, ..., 2.988, 2.994, 3.0]| 
Mezi sousedními hranicemi leží jeden bucket. 
Bucket č. $i$ pokrývá interval $[b_i, b_{i+1})$ a má šířku $w_i = b_{i+1} - b_i$.
V naší konfiguraci jsou borders typicky rovnoměrně rozložené, takže všechny buckety mají stejnou šířku. 
Implementace ale podporuje i nerovnoměrné dělení (hustší buckety v zajímavých oblastech, řidší na krajích).
\par   

Vraťme se zpátky k dekódování.
Jednotlivé prvky v schématu, který jsme uvedli na začátku sekce, plní následující funkci:
\begin{enumerate}
    \item \lstinline|Linear(512 -> 1024)|: Vezme $512$-dimenzionální embedding a promítne ho do $1024$ dimenzí. 
    To je lineární transformace — násobení maticí vah $W_1 \in \mathbb{R}^{512 \times 1024}$.
    \item \lstinline|GELU()|: Aplikuje nelineární aktivaci. 
    Bez ní by dvě lineární vrstvy za sebou byly ekvivalentní jedné lineární vrstvě (součin matic je matice). 
    GELU umožňuje decoderu zachytit nelineární vztahy mezi embeddingem a výstupní distribucí.
    \item \lstinline|Linear(1024 -> num_bars)|: Zúží z $1024$ dimenzí na \lstinline|num_bars| (řekněme $1000$). 
    Výstupem je $1000$ čísel, říkáme jim \textbf{logity}, jeden pro každý bucket.
\end{enumerate}

Logit v tomto případě je jenom surové číslo z poslední lineární vrstvy našeho dekóderu.
Může to být jakékoliv reálné číslo.
Důležité ale je, že nemají žádné jednotky, nesčítají se na jedničku a ne této etapě nemají žádnou přímou pravděpodobnostní interpretaci.
Logit je pouze signál pro určitý bucket, do kterého spadne, jak moc si model myslí, že skutečná hosnota padne do daného intervalu.
Čím vyšší je logit, tím víc si model jistý.
Ale pro lepší interpretovatelnost musí tyto hodnoty nejdříve projít softmax vrstvou.
Uvažujeme–li, že máme $1000$ bucketů, softmax převede $1000$ logitů na $1000$ pravděpodobností:
$$p_i = \frac{e^{z_i}}{\sum_{j=1}^{1000} e^{z_j}}$$
Logicky pak součet všech pravděpodobností je přesně $1$ a vyšší logit implikuje vyšší pravděpodobnost.
Hustotu pravděpodobnosti potom definujeme jako $f_i = \frac{p_i}{w_i}$, kde $w_i$ je šířka bucketu.
Širší bucket musí mít nižší hustotu při stejné pravděpodobnosti, aby distribuce dávala smysl jako aproximace spojitého rozdělení. 
Při rovnoměrných bucketech je tento rozdíl jen konstantní škálování.
Evidentně pak můžeme dopočítat i všechny potřebné diskriptivní charakteristiky jako stření hodnotu, rozptyl atd.
Jak jsme již mínili dříve, neklademe na BarDistribution žádná omezení co se tyče tvaru výsledné distribuce.
Avšak v kontextu GP inference výstup typicky vypadá přibližně jako gaussovský zvon. 
To proto, že GP posteriorní prediktivní distribuce je přesně normální rozdělení.
Pokud PFN dobře aproximuje GP posterior, buckety budou mít tvar zvonu. 
Ale model to nikdo nenutí, naučil se to sám, protože loss funkce, se kterou byl model trénován, ho odměňuje za to, že jeho distribuce co nejlépe odpovídá skutečné posteriorní distribuci.
\par   

Dále poskytneme krátký komentář, proč během trénování PFN volíme pro loss funkci právě Negative Log–Likelihood (NLL) tvar.
NLL je ztrátová funkce, která měří kvalitu predikované distribuce. Postupuje následovně:
\begin{enumerate}
    \item Model vyprodukuje logity $z_1, \dots, z_{1000}$.
    \item Softmax je převede na pravděpodobnosti $p_1, \dots, p_{1000}$.
    \item Pravděpodobnosti se převedou na hustoty: $f_i = p_i / w_i$.
    \item Najde se bucket $k$, do kterého padne skutečné $y_{\text{test}}$.
    \item Loss pro tento bod je: $\mathcal{L} = -\log(f_k) = -\log(p_k) + \log(w_k)$
\end{enumerate}

Člen $\log(w_k)$ je konstanta (borders se nemění), takže gradient teče jen přes $-\log(p_k)$. 
Minimalizace NLL je ekvivalentní maximalizaci likelihood, to znamená, že model se učí přiřadit co nejvyšší hustotu tam, kde skutečné $y$ leží.

Otázkou je, proč bychom nemohli volit např. tvar Mean Square Error (MSE) jako loss funkci.
MSE loss by měřil jen $(y_{\text{pred}} - y_{\text{true}})^2$, tj. chybu bodového odhadu. 
NLL přes BarDistribution měří kvalitu celé distribuce.
Model se učí nejen správný střed, ale i správnou šířku a tvar nejistoty. 
To je přesně to, co potřebujeme pro aproximaci GP posterioru, chceme celou PDD, ne pouze průměr.
Navíc, jak se ukáže dále (REFERENCEnagler), minimalizace NLL vede k optimálnímu řešení, které je Bayesovská posteriorní prediktivní distribuce, tedy přesně to, co GP analyticky počítá.
\par    

Na konci sekce zkusíme shrnout kompletní postup.
Na úplném začátku jsme modelu předložili sadu trénovacích a testovacích bodů, přitom ale model nedostal \texttt{y\_test} a účelem bylo ho naučit tyto hodnoty predikovat.
PFN postupně zpracovalo a přetransformovalo skrz svoje vrstvy \texttt{x\_train}, \texttt{y\_train} a \texttt{x\_test}.
To znamená, že model zpracoval mu předložená informace jako celek naraz.
Během toho se PFN naučilo pravidlo, jak se chovají různé GP posteriory pro různé parametry.
Na konci jsme dostali pravděpodobnostní rozdělení pomocí převodu na logity a použití BarDistribution.
Výsledek jsme vždy porovnali se skutečným GP a jako metriku, jak moc dobře PFN odhaduje skutečnost, jsme použili NLL, který bychom chtěli minimalizovat.
Po dostatečném tréninku jsme pak dostali model, který na základě předložených dat ze skutečneho GP prioru, které nikdy neviděl, je schopen odhadnout, odkud pochází a jak vypadá funkce, ze které pochází.


